%% ----------------------------------------------------------------------------
%%  §11.2  Proprietà delle serie
%% ----------------------------------------------------------------------------
\section{Proprietà delle serie}
\label{sec:proprieta-serie}

Poiché il carattere e la somma di una serie sono definiti mediante il limite delle somme
parziali, le proprietà dei limiti di successioni si traducono in proprietà delle serie.
Cominciamo dalla linearità.

\begin{proposizione}[Linearità]\label{prop:linearita-serie}
\index{serie!linearità}
\begin{enumerate}[label=\textup{(\roman*)}]
  \item Se le serie $\sum a_{k}$ e $\sum b_{k}$ sono regolari e la somma
  $\sum_{k=1}^{\infty}a_{k}+\sum_{k=1}^{\infty}b_{k}$ ha significato in $\overline\R$ (cioè non è della forma
  $+\infty-\infty$), allora la serie $\sum(a_{k}+b_{k})$ è regolare e
  \begin{equation}\label{eq:serie-somma}
    \sum_{k=1}^{\infty}(a_{k}+b_{k})=\sum_{k=1}^{\infty}a_{k}+\sum_{k=1}^{\infty}b_{k} .
  \end{equation}
  \item Se la serie $\sum a_{k}$ è regolare e $c\neq0$ è un numero reale, anche la serie
  $\sum c\,a_{k}$ è regolare e
  \begin{equation}\label{eq:serie-costante}
    \sum_{k=1}^{\infty}c\,a_{k}=c\sum_{k=1}^{\infty}a_{k} ,
  \end{equation}
  con le usuali regole dei segni se la serie diverge. Se $\sum a_{k}$ converge, la
  \eqref{eq:serie-costante} vale anche per $c=0$.
\end{enumerate}
\end{proposizione}

\begin{schema}
\Ipotesi{(i) $\sum a_{k}$, $\sum b_{k}$ regolari, con somme non della forma $+\infty-\infty$; \
(ii) $\sum a_{k}$ regolare, $c\neq0$.}
\Tesi{$\sum(a_{k}+b_{k})=\sum a_{k}+\sum b_{k}$; \ $\sum c\,a_{k}=c\sum a_{k}$.}
\end{schema}

\begin{strategia}
Le somme parziali di $\sum(a_{k}+b_{k})$ e di $\sum c\,a_{k}$ si ottengono da quelle delle serie date
con le operazioni corrispondenti, perché per le somme \emph{finite} le proprietà
associativa, commutativa e distributiva valgono senza restrizioni. Il passaggio al limite
è allora garantito dai teoremi sulle operazioni con i limiti, che escludono proprio le
forme indeterminate $+\infty-\infty$ e $0\cdot\infty$.
\end{strategia}

\begin{dimostrazione}
Le somme parziali di $\sum(a_{k}+b_{k})$ e di $\sum c\,a_{k}$ sono, per le proprietà delle
somme finite, $\sum_{k=1}^{n}(a_{k}+b_{k})=\sum_{k=1}^{n}a_{k}+\sum_{k=1}^{n}b_{k}$ e
$\sum_{k=1}^{n}c\,a_{k}=c\sum_{k=1}^{n}a_{k}$. La tesi segue dai teoremi sulle operazioni con i
limiti, nel caso finito (\cref{teo:operazioni-limiti}) e in quello infinito
(\cref{teo:operazioni-infinite}), esclusi appunto i casi di indeterminazione.
\end{dimostrazione}

Ad esempio, dalla serie di Mengoli si ricava subito
$\sum_{k=1}^{\infty}\frac3{k(k+1)}=3$. Si osservi che la somma di due serie indeterminate può
essere convergente: $\sum(-1)^{k}$ e $\sum(-1)^{k+1}$ sono indeterminate, ma la loro somma è la
serie con tutti i termini nulli.

\subsection{Il resto di una serie}

Le prime somme parziali di una serie possono essere qualsiasi: il carattere dipende
soltanto da ciò che accade «in fondo». Per precisare questa idea, data la serie
$\sum_{k=1}^{\infty}a_{k}$ e fissato $n\in\N$, consideriamo la serie \emph{resto $n$-esimo}
\begin{equation}\label{eq:serie-resto}
  \sum_{k=n+1}^{\infty}a_{k}=a_{n+1}+a_{n+2}+a_{n+3}+\cdots ,
\end{equation}
ottenuta trascurando i primi $n$ termini.

\begin{teorema}[del resto]\label{teo:resto-serie}\index{serie!resto}
Per ogni $n\in\N$ la serie resto $n$-esimo \eqref{eq:serie-resto} ha lo stesso carattere
della serie $\sum_{k=1}^{\infty}a_{k}$. Se la serie converge e $s$ è la sua somma, detta
$R_{n}=\sum_{k=n+1}^{\infty}a_{k}$ la somma del resto, si ha
\begin{equation}\label{eq:resto-serie}
  s=s_{n}+R_{n} \qquad\text{e}\qquad \lim_{n\to+\infty}R_{n}=0 .
\end{equation}
\end{teorema}

\begin{schema}
\Ipotesi{$\sum_{k=1}^{\infty}a_{k}$ serie qualsiasi; \ $n\in\N$ fissato.}
\Tesi{$\sum_{k=n+1}^{\infty}a_{k}$ ha lo stesso carattere di $\sum_{k=1}^{\infty}a_{k}$; \ se
$\sum a_{k}=s\in\R$, allora $R_{n}=s-s_{n}\to0$.}
\end{schema}

\begin{strategia}
Le somme parziali del resto sono differenze di somme parziali della serie data:
$\sum_{k=n+1}^{m}a_{k}=s_{m}-s_{n}$. Con $n$ fissato, $s_{n}$ è una costante, e le due successioni
(nell'indice $m$) hanno lo stesso comportamento. Che il resto tenda a zero è allora la
definizione di convergenza riletta.
\end{strategia}

\begin{dimostrazione}
Fissato $n$, per ogni $m>n$ la somma parziale $(m-n)$-esima del resto è
\[
  \sum_{k=n+1}^{m}a_{k}=\sum_{k=1}^{m}a_{k}-\sum_{k=1}^{n}a_{k}=s_{m}-s_{n} .
\]
Poiché $s_{n}$ non dipende da $m$, la successione $m\mapsto s_{m}-s_{n}$ ha limite (finito o
infinito) se e solo se lo ha $s_{m}$, e i due limiti differiscono per la costante $s_{n}$.
Dunque il resto e la serie hanno lo stesso carattere, e se $s_{m}\to s$ si ha
$R_{n}=\lim_{m\to+\infty}(s_{m}-s_{n})=s-s_{n}$. Infine $R_{n}=s-s_{n}\to s-s=0$ per $n\to+\infty$.
\end{dimostrazione}

Il teorema ha una conseguenza pratica importante: \emph{modificare un numero finito di
termini non cambia il carattere di una serie} (ne cambia, in generale, la somma). Infatti
due serie che differiscono soltanto nei primi $n$ termini hanno lo stesso resto
$n$-esimo. Per questo, nei criteri di convergenza che seguono, le ipotesi potranno essere
richieste soltanto «da un certo indice in poi». Il teorema dice anche che cosa significhi
calcolare la somma di una serie con una somma parziale: l'errore commesso è esattamente
il resto $R_{n}$, e tutto il problema consiste nello \emph{stimarlo}, come abbiamo già fatto
con la formula di Taylor nel \cref{cap:taylor}.

\subsection{Il criterio di Cauchy}

Il criterio di convergenza di Cauchy per le successioni (\cref{teo:criterio-cauchy}),
applicato alle somme parziali, fornisce una condizione necessaria e sufficiente per la
convergenza di una serie che non richiede di conoscerne la somma.

\begin{teorema}[Criterio di Cauchy per le serie]\label{teo:cauchy-serie}
\index{criterio di Cauchy!per le serie}
La serie $\sum_{k=1}^{\infty}a_{k}$ è convergente se e solo se per ogni $\epsilon>0$ esiste
$\nu\in\N$ tale che
\begin{equation}\label{eq:cauchy-serie}
  \Bigl|\sum_{k=n+1}^{n+p}a_{k}\Bigr|=\abs{a_{n+1}+a_{n+2}+\dots+a_{n+p}}<\epsilon
  \qquad\perogni n>\nu,\ \perogni p\in\N .
\end{equation}
\end{teorema}

\begin{schema}
\Ipotesi{$a_{k}$ successione di numeri reali.}
\Tesi{$\sum a_{k}$ convergente $\sse$ $\perogni\epsilon>0\ \esiste\nu$:
$\bigl|\sum_{k=n+1}^{n+p}a_{k}\bigr|<\epsilon$ per ogni $n>\nu$ e ogni $p\in\N$.}
\end{schema}

\begin{strategia}
La serie converge se e solo se la successione $s_{n}$ converge, cioè, per il criterio di
Cauchy, se e solo se $s_{n}$ è di Cauchy. Resta da tradurre la condizione
$\abs{s_{m}-s_{n}}<\epsilon$ in termini degli $a_{k}$: per $m=n+p>n$, la differenza $s_{m}-s_{n}$ è
la somma dei termini da $a_{n+1}$ a $a_{n+p}$.
\end{strategia}

\begin{dimostrazione}
Per il \cref{teo:criterio-cauchy}, la serie converge se e solo se per ogni $\epsilon>0$ esiste
$\nu$ tale che $\abs{s_{m}-s_{n}}<\epsilon$ per ogni $m,n>\nu$. Poiché $\abs{s_{m}-s_{n}}=\abs{s_{n}-s_{m}}$ e
$\abs{s_{n}-s_{n}}=0$, basta considerare le coppie con $m>n$; scritto $m=n+p$ con $p\in\N$,
\[
  s_{m}-s_{n}=\sum_{k=1}^{n+p}a_{k}-\sum_{k=1}^{n}a_{k}=\sum_{k=n+1}^{n+p}a_{k} ,
\]
e la condizione diventa la \eqref{eq:cauchy-serie}.
\end{dimostrazione}

Con $p=1$ la \eqref{eq:cauchy-serie} dà $\abs{a_{n+1}}<\epsilon$ per $n>\nu$: si ritrova la condizione
necessaria del \cref{teo:condizione-necessaria}. Il criterio dice di più: non basta che
i singoli termini siano piccoli, devono essere piccoli anche tutti i «blocchi» di termini
consecutivi, comunque lunghi. Lo useremo per dimostrare il teorema sulla convergenza
assoluta (\S\ref{sec:convergenza-assoluta}) e, nel \S\ref{sec:serie-armonica}, per una terza
dimostrazione della divergenza della serie armonica.

\begin{osservazione}[Parentesi in una serie]\label{oss:parentesi}
Raggruppare con parentesi termini consecutivi di una serie regolare non ne cambia la
somma: le somme parziali della serie raggruppata sono una estratta $s_{n_{1}},s_{n_{2}},\dots$
delle somme parziali della serie data, e un'estratta ha lo stesso limite della
successione (\cref{prop:estratta-limite}). Il viceversa è falso: la serie
$(1-1)+(1-1)+(1-1)+\cdots$ ha tutti i termini nulli e somma $0$, ma togliendo le parentesi si
ottiene la serie indeterminata $1-1+1-1+\cdots$. Le parentesi si possono togliere
liberamente soltanto se si sa già che la serie senza parentesi è regolare.
\end{osservazione}

\begin{osservazione}[Termini nulli]\label{oss:termini-nulli}
Inserire o togliere termini nulli in una serie non ne cambia né il carattere né la
somma. Togliendo gli zeri, infatti, le somme parziali della nuova serie sono un'estratta
di quelle della serie data, e ne hanno lo stesso limite; viceversa, inserendo degli zeri
le somme parziali si ripetono, nello stesso ordine, e le nuove somme parziali hanno lo
stesso limite delle vecchie (se una successione tende a $s$, vi tende anche quella
ottenuta ripetendo un numero finito di volte ciascun termine, come si verifica subito
con la definizione di limite).
\end{osservazione}
